\subsection*{Archived v2.27.3 local post-release note (historical baseline memory)}



% --- Integrated v3.31.0 groundwork-complete pre-OP frontier synthesis ---

\section{Release orientation}
This document is presented as the full monolithic Companion text of the current line. Within that full monolithic presentation, the newly integrated v3.31.0 block freezes the mathematically relevant pre-OP foundation developed in the current working stage. The release claim of that integrated block is therefore intentionally limited and precise:
\begin{enumerate}[leftmargin=2em]
  \item the groundwork layer before OP 1--5 is completed as a coherent route-faithful structure;
  \item the current local state is conservatively instantiated as a frontier-certified pre-entry state;
  \item the line already supports controlled witness-based descent and lead-guided frontier resolution;
  \item explicit OP 1--5 work is intentionally deferred to the next stage.
\end{enumerate}

\section{The three-coordinate live residual interface}
\begin{definition}[Three-coordinate live residual interface]\label{def:live-residual-interface}
The current line is organized around the live residual interface
\[
\Rlive=(\rho_{\mathrm{CL}},\rho_{\Theta},\rho_{A}),
\]
where $\rho_{\mathrm{CL}}$ is the horizon-side scalar closure coordinate, while $(\rho_{\Theta},\rho_{A})$ are the boundary-side slow phase and slow envelope packet coordinates.
\end{definition}

\begin{remark}
The current open residual work is therefore carried by exactly three live coordinates: one scalar horizon-side coordinate and a two-component boundary-side packet coordinate.
\end{remark}

\begin{definition}[Current normalized state]\label{def:current-normalized-state}
Using a current admissible normalization scale triple
\[
(s_{\mathrm{CL,curr}},s_{\Theta,\mathrm{curr}},s_{A,\mathrm{curr}}),
\]
the current normalized live residual state is
\[
\Rcurrtilde=(\widetilde{\rho}_{\mathrm{CL,curr}},\widetilde{\rho}_{\Theta,\mathrm{curr}},\widetilde{\rho}_{A,\mathrm{curr}}),
\]
with the obvious channel-wise normalization.
\end{definition}

\section{Paired conservative machine and compression}
\begin{definition}[Local paired conservative machine]\label{def:paired-machine}
The local paired conservative machine consists of the paired intake, request, response, fulfillment, resolution, close-or-promote, cycle, return, and branch layers, always preserving:
\begin{enumerate}[leftmargin=2em]
  \item the inherited horizon/boundary distinction,
  \item the inherited earliest-boundary stop,
  \item certificate-neutrality.
\end{enumerate}
\end{definition}

\begin{definition}[Compact paired conservative normal form]\label{def:compact-normal-form}
The compact paired conservative normal form is the compressed ordered object
\[
\mathcal{N}^{\ast}_{\mathrm{pair}}=(\mathcal{N}_{\mathrm{hor}},\mathcal{N}_{\mathrm{bdry}}),
\]
which records only the downstream-relevant scalar horizon-side content and packet boundary-side content carried by the paired machine.
\end{definition}

\begin{theorem}[Compression and invariant control]\label{thm:compression-invariant-control}
Every finite route-faithful local paired conservative machine admits a compact paired conservative normal form that preserves the inherited horizon/boundary distinction, the inherited earliest-boundary stop, and certificate-neutrality, while remaining readable through the extracted live residual state.
\end{theorem}

\begin{proof}
The paired machine is finite at each local stage, and its admissible content separates into one horizon-side scalar channel and one boundary-side packet channel. Compression therefore produces an ordered summary object preserving exactly the route-faithful invariants built into the original machine.
\end{proof}

\section{Tool layer}
\begin{definition}[Extracted tool layer]\label{def:tool-layer}
The extracted tool layer consists of:
\begin{enumerate}[leftmargin=2em]
  \item compression-before-comparison,
  \item monotone selection on the compressed state,
  \item a route-faithful close-or-promote shortcut,
  \item scalar-packet decoupling for genuinely one-sided arguments.
\end{enumerate}
\end{definition}

\begin{theorem}[Tool extraction theorem]\label{thm:tool-extraction}
The compact paired conservative normal form supports the extracted tool layer entirely through the preserved route-faithful invariant class and the extracted three-coordinate residual state.
\end{theorem}

\begin{proof}
Each tool is defined directly on the compact normal form and depends only on route-faithful invariants and on the residual coordinates. No reopening of the full internal machine is required unless the compressed tool layer is insufficient.
\end{proof}

\section{Tool stability and normalization}
\begin{definition}[Tool-stable invariant regime]\label{def:tool-stable-invariant-regime}
A compact paired conservative normal form lies in a tool-stable invariant regime if monotone selectors are compression-stable, close-side states are route-faithfully robust, promote-side continuation is honest rather than a pure weight artifact, and the preserved route-faithful invariants remain unchanged under admissible uses of the tool layer.
\end{definition}

\begin{definition}[Canonical normalization regime]\label{def:canonical-normalization-regime}
A canonical normalization regime is an admissible class of normalization scale triples in which the selector class and close/promote classification remain unchanged under bounded side-preserving scale replacements.
\end{definition}

\begin{theorem}[Normalization theorem]\label{thm:normalization-theorem}
Every compact paired conservative normal form equipped with an admissible normalization scale triple admits a normalized live residual reading and a normalized weighted residual energy
\[
\Ew=w_{\mathrm{CL}}|\widetilde{\rho}_{\mathrm{CL}}|^2+w_{\Theta}|\widetilde{\rho}_{\Theta}|^2+w_A|\widetilde{\rho}_{A}|^2.
\]
Inside a canonical normalization regime, comparison, selection, and close/promote testing are independent of arbitrary raw scale imbalance.
\end{theorem}

\begin{proof}
Normalization is performed channel-wise inside the route-faithful split. Within the canonical normalization regime, equivalent admissible scale choices preserve the relevant comparisons and branch classifications, so the normalized state is the preferred pre-OP working surface.
\end{proof}

\section{Residual geometry}
\begin{definition}[Normalized residual geometry]\label{def:normalized-residual-geometry}
The normalized residual space carries:
\begin{enumerate}[leftmargin=2em]
  \item admissible descent sectors and strict descent sectors,
  \item trapping close basins,
  \item honest promote regions,
  \item a residual frontier separating the geometrically unresolved zone.
\end{enumerate}
\end{definition}

\begin{definition}[Residual frontier]\label{def:residual-frontier}
The residual frontier is the part of normalized residual space lying in neither a trapping close basin nor an honest promote region.
\end{definition}

\begin{theorem}[Geometric control theorem]\label{thm:geometric-control}
In a fixed canonical normalization regime, the normalized residual space admits a route-faithful geometric control structure consisting of admissible descent sectors, close basins, honest promote regions, a residual frontier, and descent selectors on finite admissible successor families.
\end{theorem}

\begin{proof}
These structures are defined entirely within the normalized residual space and remain inside the same route-faithful invariant class. Each construction therefore preserves the inherited horizon/boundary distinction, earliest-boundary stop, and certificate-neutrality.
\end{proof}

\section{Residual dynamics}
\begin{definition}[Admissible descent trajectory]\label{def:admissible-descent-trajectory}
An admissible descent trajectory is a sequence of normalized residual states in which each successor lies in the admissible descent sector of its predecessor. If strict descent is chosen whenever available, the trajectory is called strictly selected.
\end{definition}

\begin{lemma}[First no-oscillation lemma]\label{lem:first-no-oscillation}
A strictly selected admissible descent trajectory cannot exhibit a two-step oscillation between distinct states if strict descent is available at both states.
\end{lemma}

\begin{proof}
Strict descent at the first state gives strict energy decrease; strict descent at the second state gives another strict decrease. Returning to the first state would therefore force the same energy to be both strictly larger and strictly smaller than itself, a contradiction.
\end{proof}

\begin{theorem}[Dynamic control theorem]\label{thm:dynamic-control}
Admissible descent dynamics on the normalized residual space has nonincreasing normalized weighted residual energy. Under strict descent, two-step oscillation is excluded; close-side behavior is represented by trapping close basins; promote-side behavior is represented by persistent honest promote regions; and every honest promote-to-close transition must pass through the residual frontier.
\end{theorem}

\begin{proof}
The monotonicity of energy follows from admissible descent. The no-oscillation claim is Lemma~\ref{lem:first-no-oscillation}. The basin and frontier statements are the dynamic reading of the previously defined geometric structures.
\end{proof}

\section{Compressed working principles}
\begin{definition}[Compressed working principles]\label{def:compressed-working-principles}
The current groundwork is governed by four working principles:
\begin{enumerate}[leftmargin=2em]
  \item compressed descent principle,
  \item frontier-resolution principle,
  \item basin-entry criterion,
  \item promote-persistence criterion.
\end{enumerate}
\end{definition}

\begin{theorem}[Compressed working-principle theorem]\label{thm:compressed-working-principle-theorem}
In a fixed canonical normalization regime, every local situation is to be handled by first compressing the paired machine, then normalizing the residual state, then locating the state in the close basin / promote region / residual frontier trichotomy, and only afterwards attempting any OP-specific specialization.
\end{theorem}

\begin{proof}
This is the direct synthesis of compression, normalization, geometric control, and dynamic control. The current groundwork therefore reduces local pre-OP reasoning to the compressed working principles.
\end{proof}

\section{Formal OP-entry framework}
\begin{definition}[Admissible OP-entry state]\label{def:admissible-op-entry-state}
A compact paired conservative normal form is an admissible OP-entry state if it lies in a tool-stable invariant regime and a canonical normalization regime, its normalized residual state has been classified by the resolved branch principle, and in the frontier case at least one admissible descent-guided frontier-resolution step is explicitly available.
\end{definition}

\begin{theorem}[Pre-OP foundation synthesis theorem]\label{thm:pre-op-foundation-synthesis}
The local paired conservative machine admits a complete pre-OP foundational package, together with a well-defined admissible OP-entry condition. Explicit work on OP 1--5 therefore need not begin on the raw paired machine itself, but only after passage through this mathematically controlled pre-OP foundation layer.
\end{theorem}

\begin{proof}
The package is assembled from compression, invariant control, tool extraction, tool stability, normalization, residual geometry, residual dynamics, and the compressed working principles. The admissible OP-entry state is the formal transition condition from groundwork into explicit OP work.
\end{proof}

\section{Current-state instantiation}
\begin{definition}[Current frontier-certified pre-entry state]\label{def:current-frontier-preentry}
The current local state is represented by the current compact paired conservative normal form $\mathcal{N}^{\ast}_{\mathrm{curr}}$, its extracted current live residual state $\Rcurr$, its normalized state $\Rcurrtilde$, and a current admissible normalization scale triple. In the absence of an explicit current proof of close-basin or promote-region membership, the current state is read conservatively as frontier-certified pre-entry.
\end{definition}

\begin{definition}[Preferred witness class]\label{def:preferred-witness-class}
The current preferred witness class $\Wpref$ consists of strict admissible route-faithful paired witnesses that improve the horizon-side closure coordinate and provide at least one genuine packet-side improvement without destroying the current invariant class.
\end{definition}

\begin{definition}[First normalized lead witness]\label{def:first-normalized-lead-witness}
A first normalized lead witness is an element
\[
\widetilde{\mathbf{R}}_{\mathrm{lead}}\in\Wpref
\]
that minimizes the normalized weighted residual energy inside the preferred witness class and remains admissible for downstream OP-entry auditing.
\end{definition}

\begin{definition}[Lead-direction tag]\label{def:lead-direction-tag}
The current lead-direction tag records whether the dominant current gain of the lead witness is horizon-dominant, phase-dominant, envelope-dominant, or paired-mixed.
\end{definition}

\begin{definition}[Directed next witness class]\label{def:directed-next-witness-class}
If the lead witness remains frontier-classified, the directed next witness class $\Wdir$ is obtained by filtering the current witness menu according to the current lead-direction tag, so that the next successor search is directionally coherent rather than fully undirected.
\end{definition}

\begin{theorem}[Current state synthesis theorem]\label{thm:current-state-synthesis-theorem}
For the present project line, the current local state has current foundation-complete frontier status. More precisely:
\begin{enumerate}[leftmargin=2em]
  \item the pre-OP foundational package is established;
  \item the current state is conservatively read as frontier-certified pre-entry;
  \item the current state carries a current normalization scale triple and normalized residual state;
  \item the current state carries a preferred witness class, a first normalized lead witness, a lead-direction tag, and a directed next witness class.
\end{enumerate}
Consequently, the line is mathematically stabilized up to the final pre-OP transition layer.
\end{theorem}

\begin{proof}
All listed structures have already been established in the current groundwork line and then instantiated at the local level. Collecting them yields the claimed synthesis.
\end{proof}

\section{Release conclusion}
\begin{corollary}[Groundwork-complete frontier release reading]\label{cor:groundwork-complete-frontier-release-reading}
The present line is release-ready as a groundwork-complete pre-OP frontier line. The current state is not yet claimed to be full OP-entry and does not yet claim closure of OP 1--5. Instead, the release freezes the completed pre-OP foundation together with a controlled witness-based transition regime for the next OP-focused stage.
\end{corollary}

\begin{proof}
This is the direct release reading of the current-state synthesis theorem and the formal OP-entry framework.
\end{proof}

\section*{Release note}
\addcontentsline{toc}{section}{Release note}
This release should be read as the final pre-OP launch point of the current line. The broad mathematical architecture has been stabilized and frozen. The remaining task is no longer to widen the framework further, but to carry the current frontier-certified pre-entry state through the final controlled transition into immediate OP-entry and then into explicit OP 1--5 work.

\appendix
\setcounter{section}{0}
\renewcommand{\thesection}{H\arabic{section}}
\renewcommand{\thesubsection}{H\arabic{section}.\arabic{subsection}}
